\clearpage
\section{05 关系代数：中文学习笔记}\label{05_relational_algebra_zh}

CSC5041 · 第一次课堂测验复习 · 2026-10-19

\subsection{阅读说明与学习目标}\label{05_relational_algebra_zh-ux9605ux8bfbux8bf4ux660eux4e0eux5b66ux4e60ux76eeux6807}

本篇接续 \hyperref[03_relational_model_and_integrity_zh]{03 关系模型} 和 \hyperref[04_er_to_relational_mapping_zh]{04 ER 映射}，主要依据教师 \href{../../../slides/Relational_Algebra.pptx}{Relational\_Algebra.pptx}，共 61 页。下文未注明文件的页码均指该 PPT 的幻灯片顺序。

范围对应教学计划第 1 页 Week 4；Tutorial 2 第 1 页 Question 1(a)(i)--(vi) 涉及选择投影、连接、集合运算、未参加任何课程、所有目标课程以及分组聚合。聚合虽不是经典基本关系代数的一部分，但课件第 58--59 页和 Tutorial 均有覆盖，故本篇一起整理。SQL 留到 06。

学完后，应能解释每个运算的输入与输出模式，从自然语言组合表达式，并识别错误连接、过早投影、同名混淆和除法候选范围问题。自编数据和结果是静态教学推演，不是真实数据、数据库运行结果或新的模拟卷。

\subsection{1. 关系代数的记号与闭包}\label{05_relational_algebra_zh-ux5173ux7cfbux4ee3ux6570ux7684ux8bb0ux53f7ux4e0eux95edux5305}

来源：第 2--4、55--61 页。

关系代数以关系为输入，产生新的关系，结果可以继续参与运算，这称为\textbf{闭包 closure}。课件将它称为过程式 DML：表达式给出逻辑运算的组合，而不是数据库底层必须照做的物理执行顺序。

\begin{longtable}[]{@{}
  >{\raggedright\arraybackslash}p{(\linewidth - 4\tabcolsep) * \real{0.3333}}
  >{\raggedright\arraybackslash}p{(\linewidth - 4\tabcolsep) * \real{0.3333}}
  >{\raggedright\arraybackslash}p{(\linewidth - 4\tabcolsep) * \real{0.3333}}@{}}
\toprule\noalign{}
\begin{minipage}[b]{\linewidth}\raggedright
运算
\end{minipage} & \begin{minipage}[b]{\linewidth}\raggedright
本篇写法
\end{minipage} & \begin{minipage}[b]{\linewidth}\raggedright
主要作用
\end{minipage} \\
\midrule\noalign{}
\endhead
\bottomrule\noalign{}
\endlastfoot
选择 Selection & \(\sigma _{\mathsf{condition}}(R)\) & 选行 \\
投影 Projection & \(\pi _{\mathsf{attributes}}(R)\) & 选列并去重 \\
重命名 Rename & \(\rho _{\mathsf{NewName}(\mathsf{newAttributes})}(R)\) & 区分关系角色或对齐属性名 \\
并、交、差 & \(R\,\cup \,S\)、\(R\,\cap \,S\)、\texttt{R\ −\ S} & 合并、共有、左有右无 \\
笛卡尔积 Product & \texttt{R\ ×\ S} & 所有元组配对 \\
条件连接 Theta join & \(R\,\bowtie _{\mathsf{condition}}\,S\) & 满足条件的元组配对 \\
自然连接 Natural join & \(R\,\bowtie \,S\) & 所有同名属性相等，合并这些属性 \\
除法 Division & \(R\,\div \,S\) & 与 S 中每个目标都关联的候选 \\
分组聚合 Aggregation & \(\gamma _{\mathsf{groups};\,\mathsf{aggregates}}(R)\) & 每组统计 \\
\end{longtable}

本篇用代码块中的 \texttt{\{...\}} 表示下标，便于在 Markdown 中稳定阅读；\texttt{←} 表示给中间关系命名，不修改原始数据。关系/属性重命名不改变已有值，完整列重命名列表按原属性顺序一一对应。

\textbf{默认采用集合语义且例子不含 NULL}。π 会消除重复；没有隐含行顺序。SQL 的重复值、NULL 和三值逻辑不能不加说明地搬进这些经典运算。Tutorial 中含 NULL 的 grade 不影响只使用其他属性的表达式；NULL 查询语义留到 06。

\subsection{2. 贯穿例子：先写清输入}\label{05_relational_algebra_zh-ux8d2fux7a7fux4f8bux5b50ux5148ux5199ux6e05ux8f93ux5165}

以下是自编的简化选课数据库，\textbf{不是 Tutorial 原模式}：省去 term，Takes 只表示学生是否选过某门课。

\begin{StudyCode}
Student(sid, sname, major)          PK{sid}
Course(cid, title, dept, credits)   PK{cid}
Takes(sid, cid)                    PK{sid,cid}
Takes.sid → Student.sid; Takes.cid → Course.cid
\end{StudyCode}

Student：

\begin{longtable}[]{@{}lll@{}}
\toprule\noalign{}
sid & sname & major \\
\midrule\noalign{}
\endhead
\bottomrule\noalign{}
\endlastfoot
S1 & Amy & CS \\
S2 & Ben & Math \\
S3 & Amy & CS \\
S4 & Dana & Math \\
\end{longtable}

Course：

\begin{longtable}[]{@{}llll@{}}
\toprule\noalign{}
cid & title & dept & credits \\
\midrule\noalign{}
\endhead
\bottomrule\noalign{}
\endlastfoot
C1 & Databases & CS & 4 \\
C2 & AI & CS & 3 \\
C3 & Linear Algebra & Math & 4 \\
\end{longtable}

Takes：

\begin{longtable}[]{@{}ll@{}}
\toprule\noalign{}
sid & cid \\
\midrule\noalign{}
\endhead
\bottomrule\noalign{}
\endlastfoot
S1 & C1 \\
S1 & C2 \\
S2 & C1 \\
S2 & C3 \\
S3 & C3 \\
\end{longtable}

S1、S3 同名，故\textbf{姓名不是标识}。S3 的专业是 CS，却只选了 Math 课程，故\textbf{学生专业与课程院系不是同一个条件}。S4 没有选课，后面判断``从未''和零记录情况时需要保留他。

\Needspace{11cm}
\subsection{3. 选择 σ：筛选元组，保留全部属性}\label{05_relational_algebra_zh-ux9009ux62e9-ux3c3ux7b5bux9009ux5143ux7ec4ux4fddux7559ux5168ux90e8ux5c5eux6027}

\StudyVisual{operators}{zh}

来源：第 5--10、18 页。

\StudyFormula{\sigma _{\mathsf{dept}=\texttt{'CS'}\,\land \,\mathsf{credits}>3}(\mathsf{Course})}

结果模式仍是 \texttt{(\allowbreak{}cid,\allowbreak{}title,\allowbreak{}dept,\allowbreak{}credits)\allowbreak{}}，本例只含 \texttt{(\allowbreak{}C1,\allowbreak{}Databases,\allowbreak{}CS,\allowbreak{}4)\allowbreak{}}。选择不自动删列，也不修改被保留元组的值。条件可以比较属性与常量，或比较两个属性，并用 \(\land \)、\(\lor \)、\texttt{¬} 组合。

连续选择可合并，并可交换顺序：

\StudyFormula{\sigma _{p}(\sigma _{q}(R))\,=\,\sigma _{p\land q}(R)\,=\,\sigma _{q}(\sigma _{p}(R))}

前提是相关条件都对 R 的属性有定义。本例``CS 课程且超过 3 学分''可用两个连续选择实现；它不等于``CS 课程或超过 3 学分''。

\textbf{不要混淆 SQL 的 SELECT 与关系代数的 selection}：σ 选行，SQL 的 SELECT 列表主要决定输出列。

\subsection{4. 投影 π：保留指定属性并消除重复}\label{05_relational_algebra_zh-ux6295ux5f71-ux3c0ux4fddux7559ux6307ux5b9aux5c5eux6027ux5e76ux6d88ux9664ux91cdux590d}

来源：第 11--18 页。

\StudyFormula{\pi _{\mathsf{dept}}(\mathsf{Course})\,=\,\{(\mathsf{CS}),\,(\mathsf{Math})\}}

虽然有两门 CS 课程，结果也只有一个 CS，因为结果是集合。投影后的模式只有 dept，cid、title、credits 已不在结果中。

将选择与投影组合：

\StudyFormula{\pi _{\mathsf{title}}(\sigma _{\mathsf{dept}=\texttt{'CS'}\,\land \,\mathsf{credits}>3}(\mathsf{Course}))}

结果模式为 \texttt{(\allowbreak{}title)\allowbreak{}}，本例是 \texttt{\{(Databases)\}}。从括号最内层向外读：先按条件筛课程，再取标题。

\subsubsection{不能过早丢失后续需要的属性}\label{05_relational_algebra_zh-ux4e0dux80fdux8fc7ux65e9ux4e22ux5931ux540eux7eedux9700ux8981ux7684ux5c5eux6027}

\StudyFormula{\text{正确：π\_\{title\}(σ\_\{dept='CS'\}(Course))}}
\StudyFormula{\text{错误：σ\_\{dept='CS'\}(π\_\{title\}(Course))}}

第二式已经投影掉 dept，条件无法求值。一般而言，只有条件 p 用到的全部属性都保留在 X 中时，才可写：

\StudyFormula{\pi _{X}(\sigma _{p}(R))\,=\,\sigma _{p}(\pi _{X}(R))}

嵌套投影在 \(X\,\subseteq \,Y\) 时满足 \(\pi _{X}(\pi _{Y}(R))\,=\,\pi _{X}(R)\)。若外层想要内层已删除的属性，表达式没有定义，不是``结果为空''。

\subsection{5. 重命名 ρ：处理角色和同名列}\label{05_relational_algebra_zh-ux91cdux547dux540d-ux3c1ux5904ux7406ux89d2ux8272ux548cux540cux540dux5217}

来源：第 55--56、60 页。

重命名可用于给集合运算对齐模式、区分自连接的两个角色，也能避免积和条件连接中的属性冲突。例如要找专业相同的不同学生对：

\StudyFormula{A\,\leftarrow \,\rho _{A(\mathsf{aid},\mathsf{aname},\mathsf{amajor})}(\mathsf{Student})}
\StudyFormula{B\,\leftarrow \,\rho _{B(\mathsf{bid},\mathsf{bname},\mathsf{bmajor})}(\mathsf{Student})}
\StudyFormula{\mathsf{Pairs}\,\leftarrow \,\pi _{\mathsf{aid},\mathsf{bid}}(A\,\bowtie _{\mathsf{amajor}=\mathsf{bmajor}\,\land \,\mathsf{aid}<\mathsf{bid}}\,B)}

本例约定学号顺序 \texttt{S1\textless{}S2\textless{}S3\textless{}S4}，输出为 \texttt{\{(S1,S3),(S2,S4)\}}。\texttt{aid\textless{}bid} 同时排除自己和反向重复配对；若只用 \(\mathsf{aid}\ne \mathsf{bid}\)，会同时返回 \texttt{(\allowbreak{}S1,\allowbreak{}S3)\allowbreak{}} 与 \texttt{(\allowbreak{}S3,\allowbreak{}S1)\allowbreak{}}。

重命名只改变关系/属性标识，不是把 Amy 改成另一个姓名。不能把未重命名的 Student 与自身直接自然连接，当作所有不同学生配对；那会按所有同名属性相等来匹配。

\subsection{6. 并、交、差：先保证模式相容}\label{05_relational_algebra_zh-ux5e76ux4ea4ux5deeux5148ux4fddux8bc1ux6a21ux5f0fux76f8ux5bb9}

来源：第 19--32、56 页。

\begin{longtable}[]{@{}lll@{}}
\toprule\noalign{}
运算 & 保留哪些元组？ & 是否可交换？ \\
\midrule\noalign{}
\endhead
\bottomrule\noalign{}
\endlastfoot
\(R\,\cup \,S\) & 在 R 或 S 中，包含两者都有的情况 & 是 \\
\(R\,\cap \,S\) & 同时在 R 和 S 中 & 是 \\
\texttt{R\ −\ S} & 在 R 中但不在 S 中 & 通常否 \\
\end{longtable}

三者都要求 \textbf{union compatibility 并相容}。课件第 19 页采用同名且同域的一一对应属性要求；因此应先投影出相同属性组，必要时重命名并对齐。列数一样但含义/域不相容，不能直接相减。

用本例建立两个\textbf{学号集合}：

\StudyFormula{\mathsf{CSIDs}\,\,\,\leftarrow \,\pi _{\mathsf{sid}}(\mathsf{Takes}\,\bowtie \,\sigma _{\mathsf{dept}=\texttt{'CS'}}(\mathsf{Course}))}
\StudyFormula{\mathsf{MathIDs}\,\leftarrow \,\pi _{\mathsf{sid}}(\mathsf{Takes}\,\bowtie \,\sigma _{\mathsf{dept}=\texttt{'Math'}}(\mathsf{Course}))}

得到 \texttt{CSIDs=\{S1,S2\}}，\texttt{MathIDs=\{S2,S3\}}；以下简写省去单属性元组的括号。

\begin{longtable}[]{@{}lll@{}}
\toprule\noalign{}
表达式 & 本例结果 & 自然语言 \\
\midrule\noalign{}
\endhead
\bottomrule\noalign{}
\endlastfoot
\(\mathsf{CSIDs}\,\cup \,\mathsf{MathIDs}\) & \texttt{\{S1,S2,S3\}} & 至少选过其中一类 \\
\(\mathsf{CSIDs}\,\cap \,\mathsf{MathIDs}\) & \texttt{\{S2\}} & 两类都选过 \\
\texttt{CSIDs\ −\ MathIDs} & \texttt{\{S1\}} & 选过 CS，但没选过 Math \\
\texttt{MathIDs\ −\ CSIDs} & \texttt{\{S3\}} & 选过 Math，但没选过 CS \\
\end{longtable}

这里的 or 为包含两者都成立的 inclusive or。``只选其中一类''需另行表达，例如 \((\mathsf{CSIDs}- \mathsf{MathIDs})\,\cup \,(\mathsf{MathIDs}- \mathsf{CSIDs})\)。

\subsubsection{``从未选课''不是一条记录上的不等条件}\label{05_relational_algebra_zh-ux4eceux672aux9009ux8bfeux4e0dux662fux4e00ux6761ux8bb0ux5f55ux4e0aux7684ux4e0dux7b49ux6761ux4ef6}

\StudyFormula{\mathsf{NeverIDs}\,\leftarrow \,\pi _{\mathsf{sid}}(\mathsf{Student})\,- \,\pi _{\mathsf{sid}}(\mathsf{Takes})}

本例是 \texttt{\{S4\}}。先给出全部候选学生，再减去至少有一条选课记录的学生。不应以某条记录满足 \(\mathsf{cid}\ne \mathsf{C1}\) 来表示``从未选 C1''：同一学生可能既选 C1，也选 C3。

\subsection{7. 积与连接：从所有配对到正确配对}\label{05_relational_algebra_zh-ux79efux4e0eux8fdeux63a5ux4eceux6240ux6709ux914dux5bf9ux5230ux6b63ux786eux914dux5bf9}

来源：第 33--47、61 页。

\subsubsection{笛卡尔积 ×}\label{05_relational_algebra_zh-ux7b1bux5361ux5c14ux79ef}

R 的每一行与 S 的每一行配对。若 R 有 m 行、S 有 n 行，则积有 \texttt{m×n} 行；在属性名已消歧的前提下，度为两者度之和。本例 Student×Course 有 \texttt{4×3=12} 行、7 个属性。

积不会自动按主外键匹配。如果有同名属性，先用 ρ 消歧；本篇在条件连接中也用 \texttt{关系名.\allowbreak{}属性名} 指明来源，可理解为对冲突属性作了相应限定重命名。

\subsubsection{θ 连接与等值连接}\label{05_relational_algebra_zh-ux3b8-ux8fdeux63a5ux4e0eux7b49ux503cux8fdeux63a5}

\StudyFormula{R\,\bowtie _\theta \,S\,=\,\sigma _\theta (R\,\times \,S)}

该等式在属性名已消歧、条件有定义时成立。θ 可以是比较条件；若连接条件只使用相等比较，就是等值连接 equi-join。

\StudyFormula{\mathsf{Takes}\,\bowtie _{\mathsf{Takes}.\mathsf{cid}=\mathsf{Course}.\mathsf{cid}}\,\mathsf{Course}}

这只配对课程编号相同的行。本例有 5 行。普通等值连接保留两边的连接属性副本（可用限定名区分）；\textbf{不会因两列值相等就自动把它们合并为一列}。

\subsubsection{自然连接 ⋈}\label{05_relational_algebra_zh-ux81eaux7136ux8fdeux63a5}

自然连接强制\textbf{所有同名属性}相等，并只保留这些属性的一份。不存在同名属性时，它退化为笛卡尔积。

\StudyFormula{\mathsf{Student}\,\bowtie \,\mathsf{Takes}\,\bowtie \,\mathsf{Course}}

在本例中，第一步只有 sid 同名，接着只有 cid 同名；结果模式为 \texttt{(\allowbreak{}sid,\allowbreak{}sname,\allowbreak{}major,\allowbreak{}cid,\allowbreak{}title,\allowbreak{}dept,\allowbreak{}credits)\allowbreak{}}，共 5 行。

自然连接不会理解``哪些列在业务上应该匹配''。若 Instructor 与 Course 都有 dept，自然连接可能额外要求教师院系等于课程院系；若两者不必相同，就应采用明确连接条件或先重命名无关的同名属性。

第 37 页把等值筛选后去掉一份连接列的结果联系到自然连接；严格使用时仍要按第 40--41 页检查所有同名属性，不能把任意``等值连接后投影''一概称为自然连接。连接只保留有匹配的组合，不自动给没有选课的 S4 补一行。

\subsection{8. 组合查询：保留身份，最后输出姓名}\label{05_relational_algebra_zh-ux7ec4ux5408ux67e5ux8be2ux4fddux7559ux8eabux4efdux6700ux540eux8f93ux51faux59d3ux540d}

来源：第 25--32、43--47 页；以下同名反例为自编。

找``既选过 CS 又选过 Math 的学生姓名''：先求学号交集，再连接 Student：

\StudyFormula{\mathsf{BothIDs}\,\leftarrow \,\mathsf{CSIDs}\,\cap \,\mathsf{MathIDs}}
\StudyFormula{\mathsf{Answer}\,\,\leftarrow \,\pi _{\mathsf{sname}}(\mathsf{Student}\,\bowtie \,\mathsf{BothIDs})}

本例返回 \texttt{\{(Ben)\}}。若先分别取姓名，两边都含 Amy：CS 来自 S1，Math 来自 S3。姓名交集会错误返回 Amy，即把\textbf{两个不同的人}合成一个满足两条件的人。

对``选过 CS 但没选过 Math''也应先作学号差。本例正确者 S1 名叫 Amy，但两边姓名集相减会把 Amy 删掉。\textbf{投影对并可分配，但一般不能对交和差随意提前投影。}

输出只要姓名时，最后的 π 会合并同名结果，这是题目要求的姓名集合；若需要区分学生，应保留 sid。不要用输出姓名这一要求反推姓名是键。

\subsubsection{Tutorial 的 term 不能随意丢掉}\label{05_relational_algebra_zh-tutorial-ux7684-term-ux4e0dux80fdux968fux610fux4e22ux6389}

Tutorial 2 的 Enrolled 和 Teaches 都包含 cid、term。``上过某教师教的课程''需要同一课程、同一开课 term 对应；只按 cid 连接，可能把某学生在另一学期的选课错配给该教师。

本篇 Takes 特意省去 term，不能直接代替 Tutorial 的完整输入模式；套用模板前必须先核对关联所需的完整属性组。

\Needspace{11cm}
\subsection{9. 除法 ÷：表达``每一个目标都满足''}\label{05_relational_algebra_zh-ux9664ux6cd5-ux8868ux8fbeux6bcfux4e00ux4e2aux76eeux6807ux90fdux6ee1ux8db3}

\StudyVisual{division}{zh}

来源：第 48--54、61 页。

设 R 的属性为不相交的 X 与 Y，S 的属性恰为 Y。\(R\,\div \,S\) 输出 X，保留那些在 R 中与 \textbf{S 的每个元组}都配对的 X 值，候选取自 \(\pi _{X}(R)\)。

\StudyFormula{R(X,Y)\,\div \,S(Y)\,\rightarrow \,\mathsf{Result}(X)}

所以除法不是数值相除，也不是``匹配其中任意一个''。需要先确定候选是什么 X、要求覆盖的目标是什么 Y。

\subsubsection{本例：选过所有 CS 课程的学生}\label{05_relational_algebra_zh-ux672cux4f8bux9009ux8fc7ux6240ux6709-cs-ux8bfeux7a0bux7684ux5b66ux751f}

\StudyFormula{R\,\leftarrow \,\pi _{\mathsf{sid},\mathsf{cid}}(\mathsf{Takes})}
\StudyFormula{D\,\leftarrow \,\pi _{\mathsf{cid}}(\sigma _{\mathsf{dept}=\texttt{'CS'}}(\mathsf{Course}))}
\StudyFormula{\mathsf{AllCSIDs}\,\leftarrow \,R\,\div \,D}

D 为 \texttt{\{C1,C2\}}，S1 同时关联二者，S2 只关联 C1，S3 只关联 C3，故结果为 \texttt{\{S1\}}。最后可用 \(\pi _{\mathsf{sname}}(\mathsf{Student}\,\bowtie \,\mathsf{AllCSIDs})\) 返回姓名。

\begin{itemize}
\tightlist
\item
  ``所有 CS 课程''要求集合应来自完整 Course，不应只从已有选课中提取，否则无人选过的必需课程会被漏掉。
\item
  ``所有 CS 课程''允许额外选 Math 课程，不等于``只选 CS 课程''。
\item
  被除关系只保留 \texttt{(\allowbreak{}sid,\allowbreak{}cid)\allowbreak{}}，不夹带 sname、term、grade 等无关属性，否则除法会保留它们并改变``谁需要覆盖所有目标''的含义。
\end{itemize}

\subsubsection{多属性目标与结果模式}\label{05_relational_algebra_zh-ux591aux5c5eux6027ux76eeux6807ux4e0eux7ed3ux679cux6a21ux5f0f}

课件第 53--54 页有 R(A,B,C)、S(B,C)：

\StudyFormula{\text{R ÷ S          → 输出 A，要求覆盖 S 中全部 (B,C) 配对}}
\StudyFormula{\text{R ÷ π\_C(S)     → 输出 (A,B)，要求每个固定 (A,B) 覆盖全部 C}}

两者的问题不同。按该页给定实例，第一式得到 \texttt{\{a1\}}；第二式得到 \texttt{\{(a1,b1),(a3,b2)\}}。不要只看到 ``all'' 就忽略目标关系的模式。

\subsection{10. 除法的差集展开与空目标}\label{05_relational_algebra_zh-ux9664ux6cd5ux7684ux5deeux96c6ux5c55ux5f00ux4e0eux7a7aux76eeux6807}

来源：第 57 页；候选全集与空目标的讨论是定义边界的解释性补充。

把``每个目标都有''改写成``没有缺失的目标''：

\StudyFormula{\mathsf{Candidates}\,\leftarrow \,\pi _{X}(R)}
\StudyFormula{\mathsf{Required}\,\,\,\leftarrow \,\mathsf{Candidates}\,\times \,S}
\StudyFormula{\mathsf{Missing}\,\,\,\,\leftarrow \,\mathsf{Required}\,- \,R}
\StudyFormula{\mathsf{Bad}\,\,\,\,\,\,\,\,\leftarrow \,\pi _{X}(\mathsf{Missing})}
\StudyFormula{\mathsf{Result}\,\,\,\,\,\leftarrow \,\mathsf{Candidates}\,- \,\mathsf{Bad}}

这里 R 恰为 \texttt{(\allowbreak{}X,\allowbreak{}Y)\allowbreak{}}，Required 与 R 的模式需对齐。若原输入含更多属性，先投影到 X、Y。式中的 \texttt{X,\allowbreak{}Y} 是属性组，不是拿两个关系名作集合差。

在本例中，Candidates 为 \texttt{\{S1,S2,S3\}}。Required 有 6 个学号--课程配对；缺失 \texttt{(\allowbreak{}S2,\allowbreak{}C2)\allowbreak{}}、\texttt{(\allowbreak{}S3,\allowbreak{}C1)\allowbreak{}}、\texttt{(\allowbreak{}S3,\allowbreak{}C2)\allowbreak{}}，故 Bad 为 \texttt{\{S2,S3\}}，剩下 S1。R 中的额外 Math 配对不会造成缺失，因此不影响 S1 的资格。

\subsubsection{目标为空时，必须说清候选范围}\label{05_relational_algebra_zh-ux76eeux6807ux4e3aux7a7aux65f6ux5fc5ux987bux8bf4ux6e05ux5019ux9009ux8303ux56f4}

按上式，若 S 为空，Required 和 Missing 都为空，标准除法返回 \(\pi _{X}(R)\)。它\textbf{不会自动把从未在 R 中出现的学生补进来}。

如果题意是``所有 Student 都作为候选；每人满足每个目标''，应使用明确全集 U：

\StudyFormula{U\,\,\,\,\,\,\leftarrow \,\pi _{\mathsf{sid}}(\mathsf{Student})}
\StudyFormula{\mathsf{Bad}\,\,\,\,\leftarrow \,\pi _{\mathsf{sid}}((U\,\times \,D)\,- \,R)}
\StudyFormula{\mathsf{Result}\,\leftarrow \,U\,- \,\mathsf{Bad}}

这里 R、D 与第 9 节相同。D 非空时，此式与标准除法的结果一致（在本例外键有效的前提下）；D 为空时返回全部 U，包含没有选课的 S4。这体现空集上的全称条件为真，不代表任意实际题目都默认允许无目标；若题目排除这种情形，应说明假设。

\subsection{11. γ：分组与聚合}\label{05_relational_algebra_zh-ux3b3ux5206ux7ec4ux4e0eux805aux5408}

来源：第 58--59 页；Tutorial 2 Question 1(a)(vi)。以下统一记号是本笔记的书写约定。

γ 可以配合 COUNT、SUM、AVG、MIN、MAX。\(\gamma _{\text{分}\text{组}\text{属}\text{性};\,\text{聚}\text{合}\text{函}\text{数}\rightarrow \text{结}\text{果}\text{名}}(R)\) 表示每组产生一行，保留分组属性并新增聚合值。

\StudyFormula{\mathsf{Stats}\,\leftarrow \,\gamma _{\mathsf{dept};\,\operatorname{COUNT}(\mathsf{cid})\rightarrow \mathsf{nCourses},\,\operatorname{AVG}(\mathsf{credits})\rightarrow \mathsf{avgCredits}}(\mathsf{Course})}

\begin{longtable}[]{@{}lll@{}}
\toprule\noalign{}
dept & nCourses & avgCredits \\
\midrule\noalign{}
\endhead
\bottomrule\noalign{}
\endlastfoot
CS & 2 & 3.5 \\
Math & 1 & 4 \\
\end{longtable}

再筛选课程数超过 1 的组：\(\sigma _{\mathsf{nCourses}>1}(\mathsf{Stats})\)，本例只剩 CS。聚合结果名要先在 γ 中定义，外层 σ 才能使用。

\subsubsection{数的是哪一种对象？}\label{05_relational_algebra_zh-ux6570ux7684ux662fux54eaux4e00ux79cdux5bf9ux8c61}

\begin{itemize}
\tightlist
\item
  本例 Course 一行一门课且 cid 唯一，因此 COUNT(cid) 数课程。
\item
  若先与 Takes 连接，同一课程可能按不同学生出现多行，直接计数可能数成选课配对，而非不同课程。
\item
  为数不同对象，应先在正确标识层面去重。例如从连接结果先投影 \texttt{(\allowbreak{}dept,\allowbreak{}cid,\allowbreak{}credits)\allowbreak{}}，再分组；不能仅投影 \texttt{(\allowbreak{}dept,\allowbreak{}credits)\allowbreak{}}，因为同院系两门相同学分的课可能被合并。
\item
  \(\gamma _{\mathsf{sid};\,\operatorname{COUNT}(\mathsf{cid})\rightarrow \mathsf{nTaken}}(\mathsf{Takes})\) 只产生 S1、S2、S3 的组，计数分别 2、2、1；不会自动产生 S4 的零计数组。
\end{itemize}

本例每个被聚合组非空且无 NULL。课件第 58 页提出 NULL 问题但未在此完整定义其语义；空输入或 NULL 的处理需明确约定，SQL 中的 COUNT、NULL 等细节在 06 处理。

\subsection{12. 等价改写：先检查前提}\label{05_relational_algebra_zh-ux7b49ux4ef7ux6539ux5199ux5148ux68c0ux67e5ux524dux63d0}

来源：第 10、15--17、57、60 页；下表条件为相应定义的解释。

\begin{longtable}[]{@{}ll@{}}
\toprule\noalign{}
改写 & 成立前提或含义 \\
\midrule\noalign{}
\endhead
\bottomrule\noalign{}
\endlastfoot
\(\sigma _{p}(\sigma _{q}(R))\,=\,\sigma _{p\land q}(R)\) & 条件均有定义 \\
\(\pi _{X}(\pi _{Y}(R))\,=\,\pi _{X}(R)\) & \(X\,\subseteq \,Y\) \\
\(\pi _{X}(\sigma _{p}(R))\,=\,\sigma _{p}(\pi _{X}(R))\) & p 用到的属性全部在 X 中 \\
\(R\,\cap \,S\,=\,R\,- \,(R\,- \,S)\) & R、S 并相容 \\
\(R\,\bowtie _\theta \,S\,=\,\sigma _\theta (R\times S)\) & 属性已消歧，θ 有定义 \\
\(\pi _{X}(R\cup S)\,=\,\pi _{X}(R)\cup \pi _{X}(S)\) & 输入并相容且 X 可用，采用集合语义 \\
\end{longtable}

一般不能把最后一式的 ∪ 直接换成 ∩ 或 −：投影会丢掉身份信息，第 8 节同名学生就是反例。连接中的提前投影也必须保留输出列、连接列、之后筛选所需列。

课件第 57 页给出基本运算 \(\sigma \text{、}\pi \text{、}\cup \text{、}- \text{、}\times \text{、}\rho \)，连接、交、除法可由它们导出。第 58 页的聚合则增加统计能力，\textbf{不是仅由这些经典基本运算直接获得的一般计数能力}；不要把``可导出的运算''与``加入更强表达能力的聚合''混为一谈。

\subsection{13. 从英文需求写表达式的步骤}\label{05_relational_algebra_zh-ux4eceux82f1ux6587ux9700ux6c42ux5199ux8868ux8fbeux5f0fux7684ux6b65ux9aa4}

来源：第 25--32、43--54、60--61 页；下列为整理后的解题流程。

\begin{enumerate}
\def\labelenumi{\arabic{enumi}.}
\tightlist
\item
  \textbf{先写输出模式}：要 sid、姓名、课程标题，还是每院系一行的统计？
\item
  \textbf{确定对象标识与候选范围}：同名对象用键区分；``未参加''需要完整候选全集。
\item
  \textbf{列完整连接条件}：根据题面核对主外键及 term 等上下文；不要只凭同名就自然连接。
\item
  \textbf{翻译量词与逻辑词}：at least one 通常用连接投影；or 用并；both 用交；never 用全集减存在；every 用除法或``减去缺失者''。
\item
  \textbf{保留中间需要的属性}：先过滤/连接/集合判断身份，最后投影输出展示属性。
\item
  \textbf{逐步写中间模式}：每个 σ 的条件列仍在吗？差集两边同模式吗？除法输出是不是所需候选？
\item
  \textbf{用小实例和反例检查}：同名、多个匹配、无匹配、额外课程、缺一门目标、目标为空。
\end{enumerate}

Tutorial 2(a) 的六类题分别落在：选择投影、跨关系连接、并交、差集、除法、γ。这里给的是知识与方法；本篇没有逐题复写 Tutorial 的整套答案，也没有提前生成英文模拟卷。

\subsection{14. 易混点、复习清单与英文表达}\label{05_relational_algebra_zh-ux6613ux6df7ux70b9ux590dux4e60ux6e05ux5355ux4e0eux82f1ux6587ux8868ux8fbe}

\begin{longtable}[]{@{}
  >{\raggedright\arraybackslash}p{(\linewidth - 2\tabcolsep) * \real{0.5000}}
  >{\raggedright\arraybackslash}p{(\linewidth - 2\tabcolsep) * \real{0.5000}}@{}}
\toprule\noalign{}
\begin{minipage}[b]{\linewidth}\raggedright
容易混淆
\end{minipage} & \begin{minipage}[b]{\linewidth}\raggedright
正确区分
\end{minipage} \\
\midrule\noalign{}
\endhead
\bottomrule\noalign{}
\endlastfoot
σ / π & 选行保留模式 / 选列并去重 \\
无法求值 / 空结果 & 引用不存在属性属于表达式错误 / 合法表达式可能无匹配 \\
并相容 / 列数相同 & 属性对应、名称和域按课件要求对齐 / 只有列数不足 \\
or / only one of & 包含同时满足 / 需要排除同时满足 \\
等值连接 / 自然连接 & 明确相等条件且保留副本 / 所有同名属性相等且合并列 \\
学生姓名 / 学生标识 & 可同名 / 用键追踪同一个对象 \\
``有条记录不是'' / ``从未'' & 存在一个不满足的记录 / 不存在任何满足的记录 \\
every / only & 覆盖每个目标 / 排除目标之外的对象，两者不同 \\
除法目标 / 除法候选 & S 的全部元组 / 通常从 \(\pi _{X}(R)\) 取得 \\
COUNT 连接行 / COUNT 不同课程 & 配对次数 / 需要先保留课程标识并去重 \\
\end{longtable}

以下是复习目标，不是正式模拟测试或官方分值：

\begin{itemize}
\tightlist
\item[$\square$]
  写出每个运算的输入条件、输出模式与去重行为。
\item[$\square$]
  对本例手算选择、投影、学号并交差与自然连接。
\item[$\square$]
  正确使用重命名表达同表的两个角色。
\item[$\square$]
  用同名学生反例说明为何不能过早投影姓名。
\item[$\square$]
  解释 Tutorial 连接为何需要检查 term。
\item[$\square$]
  明确除法的 X、Y 和完整目标关系，并手算结果。
\item[$\square$]
  将除法展开成缺失配对的差集运算。
\item[$\square$]
  区分空目标下的标准除法与显式全集版本。
\item[$\square$]
  计算分组统计，识别连接重复和缺失零计数组。
\item[$\square$]
  对每条等价改写说明成立前提。
\end{itemize}

可用于英文简答的表达：

\begin{quote}
Selection filters tuples while preserving the relation's attributes.
\end{quote}

\begin{quote}
Projection retains the specified attributes and eliminates duplicate tuples.
\end{quote}

\begin{quote}
Union, intersection, and difference require union-compatible relations.
\end{quote}

\begin{quote}
A natural join matches all common attributes and retains one copy of each common attribute.
\end{quote}

\begin{quote}
Perform identity-based set operations before projecting names, because different entities may share a name.
\end{quote}

\begin{quote}
Division returns candidates associated with every tuple in the divisor.
\end{quote}

\begin{quote}
Universal conditions can be expressed by removing candidates that have a missing required pairing.
\end{quote}

\begin{quote}
The object counted depends on the input tuples and the identifiers retained before aggregation.
\end{quote}

\subsection{来源与整理记录}\label{05_relational_algebra_zh-ux6765ux6e90ux4e0eux6574ux7406ux8bb0ux5f55}

\begin{itemize}
\tightlist
\item
  教师主课件：\href{../../../slides/Relational_Algebra.pptx}{Relational\_Algebra.pptx}，第 1--61 页文字已读；重点查看第 7、10、17、19、37--38、40、42、48、50、53--54、57、60--61 页原图/公式，核对选择投影条件、并相容、连接、除法与基本运算记号。
\item
  范围依据：\href{../../../syllabus/CSC5041_Teaching_Plan.pdf}{CSC5041\_Teaching\_Plan.pdf}，第 1 页 Week 4、第 2 页 Tutorial 2；参考题面：\href{../../../slides/Tutorial\%202.pdf}{Tutorial 2.pdf}，第 1 页 Question 1(a)(i)--(vi)，已查看原页及键下划线。教学计划写 Tutorial 2 为 Week 5，题面标题写 Week 6；本篇按实际题型整理，不用这处差异推断考试改期。
\item
  自编简化模式、姓名反例、空目标与候选全集说明、聚合书写约定及计数边界均已与教师原句区分；本篇不定义 SQL 的 NULL/重复值完整语义。
\item
  本版为 2026-10-20 的 LaTeX 重排；示例证据状态沿用原说明，未新增数据库实测。
\item
  输入版本 SHA-256：关系代数课件 \texttt{\seqsplit{dcacfada73c2f482ab3b83ac7c8312cb2176932906b6d3e6cac071d01932283d}}；Tutorial 2 \texttt{\seqsplit{72c016526c1c0ff02c10fcb5f7f964c1184999c684228671025800936a6cabac}}；syllabus 教学计划 \texttt{\seqsplit{31380a7fe6504e46541b6db0bb2f24859c69188bfeb9681dca73354ebeb68295}}。
\end{itemize}
